\section{Unknot and Hopf link}\label{sec:hopf}

\subsection{Hopf link}
We attach one probe to each component (variables $y_1$, $y_2$) and set $f_{ab} = \sum_{d | (a,b)} (\mu(d)/d)\, L_{a/d, b/d}(q^d, t^d, \lambda^d)$, where $L_{ab}$ is the coefficient of $y_1^a y_2^b$ in $\log\tilde Z$ (this is \eqref{eq:extract} for one probe per component). The one-probe computation uses a closed form for the Askey--Wilson coefficients which has not been proved; it is stated as Conjecture \ref{conj:c2j} in Appendix \ref{app:virtual}, and its consequences were checked against the direct computation at $N = 3, \dots, 8$.
% [H2][H3]
\begin{result}\label{cr:hopf}
For $a, b \le 4$ one has $f_{ab} \in \dfrac{1 - q}{1 - t}\,\mathbb{Z}[q^\pm, t^\pm, (t^N)^\pm]$; when $a, b$ are both even the denominator is $1 - t^2$. The same integrality holds in the Schur basis $f_{R,b}$ with infinitely many probes on the first component ($|R| \le 4$, $b \le 2$), whereas in a naive Macdonald basis extra denominators such as $(1 - qt)$ appear. Moreover, $f_{11} = -f_{10}\, z_q(\tlam - \tlam^{-1})$ and $f_{1b} = f_{11}\,\chi_{b-1}(t^{N-1})$ with $\chi_m(x) = x^m + x^{m-2} + \dots + x^{-m}$. At $t = q$ this agrees with $\hat g_{\yone,\yone} = z(\lambda - \lambda^{-1})$ of \cite{Marino09} up to the sign and the unknot factor.
\end{result}
% [O4]
\begin{result}\label{cr:hopfdenom}
For $a, b \le 4$ the denominators of the all-$d$ invariants $g_{ab}$ are $2(1+t)\times$ a monomial in the (even, even) sectors and monomials otherwise.
\end{result}
After refinement the K\"ahler variable is shifted to $\tlam = \lambda (q/t)^{1/2}$. This may be the open-string version of the quantum shift of the K\"ahler variable in \cite{AM22}, but the ratio of shifts does not match the closed-string one. For closed strings the ratio of charges of frozen and oriented states is $x^2/Q_0 = \qr^{-1}$ (conventions of \cite{AgSc12}, Appendix \ref{app:closed}), whereas for open strings the ratio of the crosscap variable $\tlam^2$ to the cover variable $\ba_{\mathrm{KN}}^2 = \lambda^2\qr^{-1}$ is $\qr^{3}$. Since the correspondence of open and closed variables ($\lambda^2$ or $\lambda^{-2}$ with $Q$) and the open-string normalization are not yet fixed, we leave this comparison open.

\subsection{Crosscap part}
On the cover we use the Hopf link of $\U(N_A)$ refined Chern--Simons theory in composite representations: $\tfrac12 f^{\cov}_{10} = -(\lambda - \lambda^{-1})/z_t$ and $\tfrac12 f^{\cov}_{11} = (\lambda - \lambda^{-1})(\tlam - \tlam^{-1})z_q/z_t$.
% [H4][H5]
\begin{result}\label{cr:hopfg}
For the crosscap part of the Hopf link, $g_{1b} = z_q(\tlam - \tlam^{-1})\,\chi_{b-1}(\tlam/q^{1/2})$ for $b \le 4$, so that $t$ enters only through $\tlam$. The hatted invariant $\hat g_{R,\yone}$ ($|R| \le 4$) is an integer Laurent polynomial: it is $(-1)^d z_q(\tlam - \tlam^{-1})\tlam^{2d-n+1}$ for hooks $R = (n-d, 1^d)$ and $0$ otherwise. At $t = q$ this agrees with $\hat g_{\ytwo,\yone} = z(\lambda^2 - 1)$, $\hat g_{\yoneone,\yone} = z(\lambda^{-2} - 1)$ of \cite{Marino09} up to $R \leftrightarrow R^T$ and an overall sign.
\end{result}

\subsection{Determination of \texorpdfstring{$E_{22}$ and $E_{24}$}{E22 and E24}}\label{sec:E22}
When the probes of both components are split into an $x$ side (brane) and a $1/x$ side (image), the sector $(a,b) = (2,2)$ has a fixed subsector with one brane box and one image box on each component, hence $n = 2$ pairs of boundaries in the sense of Proposition \ref{prop:boundary}; similarly for $(2,4)$. The index of the cover restricted to this fixed sector is the sum $\Phifix$ over its states. Write $E = \Phifix - 2\Phifix_-$, with $\Phifix_-$ the sum over the states with $\tilde I$-eigenvalue $-1$.
\begin{hypothesis}\label{hyp:QP}
The coefficients of the monomials of $\Phifix$ are numbers of states, and $0 \le (\Phifix_-)_\omega/\Phifix_\omega \le 1$ for each monomial $\omega$.
\end{hypothesis}
For knots the signs of $\Phifix$ are mixed and Hypothesis \ref{hyp:QP} cannot be used (\S\ref{sec:EC}).
% [H6][H7]
\begin{result}\label{cr:E22}
Under Hypothesis \ref{hyp:QP}, (E1) and (E2), solved as an integer program (Appendix \ref{app:ilp}),
\begin{equation}\label{eq:E22}
E_{22} = -z_q (z_q^+)^2(\tlam - \tlam^{-1})(t q^{-1/2}\tlam + q^{1/2}\tlam^{-1})\ \ (\text{unique}),\qquad
E_{24} = E_{22}\,(\tlam^2/q + q/\tlam^2)\ \ (\text{unique}).
\end{equation}
For $(4,4)$ a solution exists but is not unique (164 of the 180 monomials of the fixed sector are free).
With many probes on the first component, $\tilde I$ acts as ``exchange of the fixed sector plus Adams operation'' and $E_{R,S} = \sum_{\rho, \sigma\vdash 2} c^R_{2;\rho}c^S_{2;\sigma}\,\varepsilon_{\rho\sigma}$, where $c^R_{2;\rho}$ are the coefficients of the Adams operation $\psi^2 s_\rho$ expanded in Schur functions $s_R$. This was checked for the ten sectors $(R, S)$ with $|R| \in \{2, 4\}$, $|S| = 2$ in which the refined part is non-zero (Result \ref{cr:hopfC}), and at $t = -1$ for all $|R| = |S| = 4$ (25 pairs): solving $\varepsilon_{\rho\sigma}(-1)$ from 4 of the 25 pairs, the remaining 21 agree.
\end{result}
Compared with $\Phifix_{22} = z_q^3(\tlam - \tlam^{-1})(tq^{-1/2}\tlam - q^{1/2}\tlam^{-1})$, $E_{22}$ replaces the difference factor $z_q^2$ by the sum factor $-(z_q^+)^2$ and flips the sign in the second factor. This is the form of a twisted trace of $\tilde I$, which we apply to knots in \S\ref{sec:EC}.

% [H8]
\begin{result}\label{cr:hopfC}
With many probes on the first component ($|R| \le 4$) and two boxes on the second ($|S| \le 2$), split $\hat C_{R,S} = M(q^{-1})^{-1}\,(g - \tfrac12E/(1+t))\,M(q^{-1})^{-1\,T}$ into the $t = q$ part (written in $\tlam$) and the refined part. Among the fourteen components with $|R| \in \{2,4\}$, $|S| = 2$, the refined part vanishes for $\hat C_{(2),(1,1)}$ and $\hat C_{(1,1),(2)}$ and for two more, leaving ten non-trivial components. In all ten the refined part is divisible by $(\tlam^2 - q)(q - t)$, which in bold variables is $\bt^2(\ba^2 - \bq^2)\cdot\bq^2(\bt^2 - 1)$, and the $t = q$ part has aligned signs in each power of $\tlam$ in all components. The quotient of the refined part is aligned in 8 of the 10 components; in the transposed pair $\hat C_{(3,1),(1,1)}$, $\hat C_{(2,1,1),(2)}$ a factor $q^{1/2}(1-t)$ enters the $\tlam^2$ block and breaks the alignment.
\end{result}

\subsection{Boundary factors}
% [O10]
\begin{result}\label{cr:boundary}
The boundary factors of Proposition \ref{prop:boundary} compare with the computations as follows.
\begin{center}
\small
\begin{tabular}{llll}
\toprule
object & $n$ & boundary factor of $E$ & boundary factor of $\Phifix$ \\
\midrule
Hopf $(2,2)$, $(2,4)$ & 2 & $z_q(z_q^+)^2$ & $z_q^3$ \\
unknot with framing $\mathsf f = 2, 4$ & 1 & $z_q^+$ & $z_q$ \\
$T(2,3)$, $T(2,5)$, $p_2$ direction & 1 & $z_q^+$ & $z_q$ \\
\bottomrule
\end{tabular}
\end{center}
For the Hopf link and the framed unknot $E$ is determined independently (Result \ref{cr:E22} and Result \ref{cr:framed}) and its boundary factor is read off. For knots $E$ is determined only modulo $2(1+t)$, so the statement is one of consistency: the integer programs of \S\ref{sec:EC} have solutions only with the boundary factor $z_q^+$, and the $E$ constructed in \S\ref{sec:principle} has this factor.
\end{result}
